%% annex/p02/propositional-metatheory-supplement.tex --- retained editorial proofs
\noindent\textit{Reader note.} These three proofs support only the separately
attributed duality, interpolation, and dual-Horn additions retained after the exact
source chapter.  The exact inherited proofs remain inline in their governed
source scopes; this supplement neither duplicates nor replaces them.

\begin{annexitem}{annex:p2.1.8}{thm:p02-propositional-duality}{Duality theorem}
\begin{proof}
Induct on the displayed NNF presentation \(\alpha\).  If \(\alpha\) is a
positive literal \(p\), then
\(v\entails p\) iff \(\bar v\not\entails p\).  If it is \(\neg p\), then
\(v\entails\neg p\) iff \(v(p)=0\), iff
\(\bar v(p)=1\), iff \(\bar v\not\entails\neg p\).  The two unit
cases follow from \(\bottom^{\mathrm d}=\topv\) and
\(\topv^{\mathrm d}=\bottom\).  For conjunction,
\[
 v\entails(\beta\land\gamma)^{\mathrm d}
 \quad\text{iff}\quad
 v\entails\beta^{\mathrm d}\text{ or }v\entails\gamma^{\mathrm d}
 \quad\text{iff}\quad
 \bar v\not\entails\beta\land\gamma,
\]
by the induction hypotheses; the disjunction case is dual.  This proves the
displayed semantic identity.

The recursive clauses also give \((\alpha^{\mathrm d})^{\mathrm d}=\alpha\)
by structural induction.  If \(\alpha\Leftrightarrow\beta\), apply the
semantic identity to \(\bar v\) for each \(v\) to obtain
\(\alpha^{\mathrm d}\Leftrightarrow\beta^{\mathrm d}\); the converse follows
from involution.  If \(\alpha\) is valid, the identity makes
\(\alpha^{\mathrm d}\) false under every \(v\).  Conversely, if the dual is
unsatisfiable, apply the identity with \(v=\bar u\) to show that every \(u\)
satisfies \(\alpha\).  Finally, the syntactic clauses swap every outer
conjunction with disjunction, every clause-disjunction with a conjunction of
the same literals, and the two empty units; hence they exchange CNF and DNF.
\end{proof}
\end{annexitem}

\begin{annexitem}{annex:p2.3.9}{thm:p02-interpolant-truth-table-correctness}{Correctness of the truth-table interpolant construction}
\begin{proof}
Let \(\theta\) be the output and let \(X\) be the common-variable set.  If a
valuation \(v\) satisfies \(\varphi\), its restriction \(v|_X\) is among the
selected assignments, so \(v\) satisfies the corresponding complete
conjunction and hence \(\theta\).  Thus
\(\varphi\entails\theta\).

Suppose instead that some valuation \(v\) satisfies \(\theta\) and falsifies
\(\psi\).  It satisfies a selected conjunction, whose assignment
\(u:X\to\{0,1\}\) therefore extends to a valuation \(w\) satisfying
\(\varphi\).  Amalgamate \(w\) on the variables occurring only in
\(\varphi\), \(v\) on those occurring only in \(\psi\), and their common
value \(u\) on \(X\).  The resulting valuation satisfies \(\varphi\) and
falsifies \(\psi\), contrary to the hypothesis.  Hence
\(\theta\entails\psi\).  By construction, every variable of \(\theta\) lies
in \(X\).  The empty selected family gives \(\bottom\), and an empty
conjunction gives \(\topv\), so the boundary cases obey the same argument.
\end{proof}
\end{annexitem}

\begin{annexitem}{annex:p2.3.12}{thm:p02-dual-horn-sat}{Dual-Horn-SAT tractability}
\begin{proof}
Complement every variable throughout the instance.  A clause with at most one
negative literal becomes a clause with at most one positive literal, and the
map is a linear-time bijection between satisfying assignments (complement the
assignment in the same way).  Apply the linear-time Horn algorithm and map its
answer back.
\end{proof}
\end{annexitem}
